\documentclass{article}
\usepackage[utf8]{inputenc}
\usepackage[margin=1.0in]{geometry}
\usepackage{amsmath}
\usepackage{amsfonts}
\usepackage{amssymb}
\usepackage{enumitem}                       % custom enum labels
\usepackage{parskip}          
\usepackage{physics}
\usepackage{geometry} 
\usepackage{esint}
\geometry{
  paperwidth=18cm,
  left=8mm,
  right=8mm,
  top=8mm,
  bottom=8mm,
}
\usepackage[framemethod=TikZ]{mdframed}     % graphics and framed envs

\renewcommand{\familydefault}{\sfdefault}      % sans serifs text
\setlength{\parindent}{0pt}                    % no paragraph indentation

% region TITLE CONSTRUCTION
\newlength\mywidth
\mywidth=\wd0
\renewcommand{\contentsname}{\hangindent=\mywidth \courseid: \coursetitle \\ \medskip \LARGE{[MY NAME HERE], \semester}}
% endregion

% region FRAMED ENVIRONMENTS
\newcounter{chapter}\setcounter{chapter}{1}
\newcounter{theo}[chapter]\setcounter{theo}{0}
\newcommand{\numTheo}{\arabic{chapter}.\arabic{theo}}
\newcommand{\mdftheo}[3]{
    \mdfsetup{
        frametitle={
            \tikz[baseline=(current bounding box.east),outer sep=0pt]
            \node[anchor=east,rectangle,fill=#3]
            {\ifstrempty{#2}{\strut #1~\numTheo}{\strut #1~\numTheo:~#2}};
        },
        innertopmargin=4pt,linecolor=#3,linewidth=2pt,
        frametitleaboveskip=\dimexpr-\ht\strutbox\relax,
        skipabove=11pt,skipbelow=0pt
    }
}
\newcommand{\mdfnontheo}[3]{
    \mdfsetup{
        frametitle={
            \tikz[baseline=(current bounding box.east),outer sep=0pt]
            \node[anchor=east,rectangle,fill=#3]
            {\ifstrempty{#2}{\strut #1}{\strut #1:~#2}};
        },
        innertopmargin=4pt,linecolor=#3,linewidth=2pt,
        frametitleaboveskip=\dimexpr-\ht\strutbox\relax,
        skipabove=11pt,skipbelow=0pt
    }
}
\newcommand{\mdfproof}[1]{
    \mdfsetup{
        skipabove=11pt,skipbelow=0pt,
        innertopmargin=4pt,innerbottommargin=4pt,
        topline=false,rightline=false,
        linecolor=#1,linewidth=2pt
    }
}
\newenvironment{theorem}[1][]{
    \refstepcounter{theo}
    \mdftheo{Theorem}{#1}{red!25}
    \begin{mdframed}[]\relax
}{\end{mdframed}}

\newenvironment{lemma}[1][]{
    \refstepcounter{theo}
    \mdftheo{Lemma}{#1}{red!15}
    \begin{mdframed}[]\relax
}{\end{mdframed}}

\newenvironment{corollary}[1][]{
    \refstepcounter{theo}
    \mdftheo{Corollary}{#1}{red!15}
    \begin{mdframed}[]\relax
}{\end{mdframed}}

\newenvironment{definition}[1][]{
    \mdfnontheo{Definition}{#1}{blue!20}
    \begin{mdframed}[]\relax
}{\end{mdframed}}

\newenvironment{example}[1][]{
    \mdfnontheo{Example}{#1}{yellow!40}
    \begin{mdframed}[]\relax
}{\end{mdframed}}

\newenvironment{proof}[1][]{
    \mdfproof{black!15}
    \begin{mdframed}[]\relax
\textit{Proof. }}{\end{mdframed}}
% endregion

% region NEW COMMANDS
\newcommand{\ds}{\displaystyle}
\newcommand{\pfn}[1]{\textrm{#1}}  % enables new functions
\newcommand{\mbf}[1]{\mathbf{#1}}  % mathbf
\newcommand{\C}{\mathbb{C}}        % fancy C
\newcommand{\R}{\mathbb{R}}        % fancy R
\newcommand{\Q}{\mathbb{Q}}        % fancy Q
\newcommand{\Z}{\mathbb{Z}}        % fancy Z
\newcommand{\N}{\mathbb{N}}   
\newcommand{\K}{\mathbb{K}}  % fancy N
\newcommand{\V}{\mathbf{V}} %vector space 
\newcommand{\0}{\mathbf{0}} %zero vector 
\newcommand{\from}{\leftarrow}
\renewcommand{\i}[1]{\textit{#1}}
\renewcommand{\b}[1]{\textbf{#1}}
\newcommand{\qed}{$\square$}
\newcommand{\range}{\text{range }}
\newcommand{\inner}[2]{\langle #1, #2 \rangle}
\title{Quantum Algorithms}
\author{Hrishikesh Belagali}
\begin{document}
\maketitle
\section{Basics}
\subsection{Basic Gates}
\begin{definition}[XOR]
    Let $a,b \in \{0,1\}^n$. The \textbf{exclusive OR} (XOR) of $a$ and $b$, denoted by $a \oplus b$, is defined as:
    \[
        (a \oplus b)_i = 
        \begin{cases}
            0 & \text{if } a_i = b_i \\
            1 & \text{if } a_i \neq b_i
        \end{cases}
    \]
    for each $i = 1, 2, \ldots, n$.
\end{definition}
\begin{definition}[FANOUT]
    Let $a \in \{0,1\}^n$ and $k \in \mathbb{N}$. The \textbf{FANOUT} of $a$ to $k$ copies, denoted by $\text{FANOUT}(a, k)$, is defined as the concatenation of $k$ copies of $a$:
    \[
        \text{FANOUT}(a, k) = \underbrace{a \| a \| \ldots \| a}_{k \text{ times}} \in \{0,1\}^{nk}
    \]
    The FANOUT is denoted by a solid dot with multiple outputs.
\end{definition}
\textbf{Note: } $\text{FANOUT}(a) = \text{CNOT}(a,0)$
When you have a CNOT with a hollow dot, it means you are copying the negation of the control qubit to 
the target qubit.
\subsection{Universal Collection of Gates}
We will first define what it means for a collection of gates to be universal for classical computing.
\begin{definition}[Universal gateset]
    A collection of logic gates is said to be universal for classical computing
    iff 
    $$\forall m, n \in \N, \forall f: \{0,1\}^m \to \{0,1\}^n,$$
    there exists a Boolean circuit constructed from the gates in the gateset that computes $f$.
\end{definition}
The following gatesets are universal for classical computing: 
\begin{itemize}
    \item $\{\text{NAND}\}$
    \item $\{\text{AND}, \text{OR}, \text{NOT}\}$ (You don't need the OR here)
    \item $\{\text{Toffoli}\}$
\end{itemize}
The proofs can be done by induction on $n$, the number of output bits.
\begin{definition}[Ancilla bits]
    The \textbf{ancilla bits} are extra bits initialized to either 0 or 1 that can be used in the computation of a Boolean function.    
\end{definition}
\subsection{Reversible Gates and Circuits}
A reversible gate/circuit is one where you can uniquely recover the input from the output. Naturally,
it would make sense for the Boolean function it implements to be reversible.
\begin{definition}[Reversible gate/circuit]
    A logic gate/circuit is said to be reversible if it computes a bijective function from 
    $\{0,1\}^n$ to $\{0,1\}^n$ for some $n \in \N$.
\end{definition}
\begin{example}
    The NOT gate is reversible and its inverse is itself.
\end{example}
\begin{example}
    The CNOT gate is reversible and its inverse is itself. This makes more sense if you write 
    it as a function:
    \[
        \text{CNOT}(a,b) = (a, a \oplus b)
    \]
    To recover $(a,b)$ from $(a, a \oplus b)$, we can do:
    \[
        (a, (a \oplus b) \oplus a) = (a,b)
    \]
\end{example}
An irreversible circuit can be made reversible through the use of ancilla bits. 
\begin{definition}[Register]
    A register is a collection of bits (or qubits).
\end{definition}
\begin{lemma}[Reversible circuit construction using Toffoli gates]
    For any $m, n \in \N$ and any $f : \{0,1\}^n \to \{0,1\}^m$, there exists
    $p \in \N$, a bit string $a \in \{0,1\}^p$ and $g : \{0,1\}^{n} \to \{0,1\}^{n+p - m}$ such that: 
    there is a reversible circuit $C$ constructed from Toffoli gates that implements 
    $\{ 0, 1 \}^n \times \{0, 1\}^p \to \{0, 1\}^m \times \{0, 1\}^{n+p - m}$ defined by:
    $(x, a) \mapsto (f(x), g(x))$. $g(x)$ is called the garbage function.
\end{lemma}
\end{document}